\subsection{Restrições Operacionais}
\label{sec:restricoes_operacionais}

As restrições operacionais definem tanto os limites mecânicos e dinâmicos do manipulador robótico quanto as condições de aproximação entre as espaçonaves.
No caso do MR, o espaço de trabalho (\emph{workspace}) é o alcance geométrico de seus elos e pela rotação permitida em cada junta.
A ponta do MR (\emph{end-effector} ou $\{E\}$) deve permanecer dentro desse espaço de trabalho, não ultrapassando os limites estabelecidos. 
Durante a fase de aproximação final, a trajetória da espaçonave visitante deve respeitar um corredor de chegada seguro (\emph{keep-out zone} ou KOZ).
Esse corredor é modelado como um cone em torno do eixo longitudinal ($R-bar$) da porta de acoplamento, restringindo a aproximação em relação ao eixo principal e também no plano lateral.
Existem também os limites dinâmicos dos atuadores. Cada junta revoluta do manipulador possui um torque máximo, enquanto a junta prismática está sujeita a uma força limite:

\begin{equation}
\label{eq:limite_torque}
|\tau_j| \leq \tau_{j,\max}, \qquad \text{juntas revolutas},
\end{equation}

\begin{equation}
\label{eq:limite_forca}
|f_5| \leq f_{5,\max}, \qquad \text{junta prismática}.
\end{equation}

Na aproximação de proximidade, adota-se também um limite para a velocidade relativa da espaçonave visitante em relação ao alvo, bem como uma distância mínima de segurança antes da entrada no corredor de acoplamento.